\chapter{Conceitos Teóricos e Contextualização do Problema}
\label{ch:fundamentos}

O Capítulo~\ref{ch:fundamentos} introduz os conceitos teóricos sobre
os quais a dissertação assenta. A Secção~\ref{sec:grafos} revisita
conceitos basilares de Teoria dos Grafos, incluindo a noção de
\textbf{densidade} (grafos esparsos e densos). A
Secção~\ref{sec:algoritmos} apresenta os algoritmos de ranqueamento
utilizados, sublinhando que todas as métricas aqui consideradas são
analisadas no contexto de \textbf{redes direcionadas}, e introduz os
coeficientes de correlação de \textbf{Spearman} e \textbf{Kendall}
que servem de critério de avaliação dos resultados. A
Secção~\ref{sec:ast} apresenta a análise estática via Árvores de
Sintaxe Abstrata (\textbf{AST}), técnica que sustenta a extração
automática do grafo a partir do código-fonte e que fundamenta a
validação do \emph{pipeline} fase a fase. A Secção~\ref{sec:re}
contextualiza o problema da priorização de requisitos, posiciona este
trabalho em relação à literatura existente --- em particular, em
relação a \citet{silva2018} --- e identifica a contribuição
incremental da presente dissertação.

\noindent\textbf{Conceitos-chave:} grafos; densidade; grafos esparsos
e densos; passeios aleatórios; PageRank; HITS; SALSA; correlação de
ranks (Spearman, Kendall); análise estática; AST; priorização de
requisitos; rastreabilidade.

% ---------------------------------------------------------------------------
\section{Conceitos básicos de Grafos}
\label{sec:grafos}

Um grafo é uma estrutura matemática que modela relações entre objetos
\parencite{newman2010, szwarcfiter2018}. Formalmente, um grafo
$G = (V,E)$ é definido por um conjunto de vértices (ou nós) $V$ e um
conjunto de arestas $E$, onde cada aresta é um par não ordenado (ou
ordenado) de vértices. Os vértices representam as entidades de
interesse, e as arestas representam as relações entre elas. Em redes
complexas, os termos ``grafo'' e ``rede'' são frequentemente usados
como sinónimos \parencite{newman2010}.

\paragraph{Exemplo.} Considere uma pequena rede de colaboração entre
investigadores. Os vértices podem ser $\{A,B,C,D\}$. As arestas
representam coautorias: se $A$ e $B$ publicaram juntos, existe uma
aresta entre eles. Suponha as colaborações $A$--$B$, $A$--$C$,
$B$--$C$, $C$--$D$. O grafo correspondente tem $V = \{A,B,C,D\}$ e
$E = \{AB,AC,BC,CD\}$.

Uma forma comum de representar um grafo é por meio da \textbf{matriz
de adjacência} $A$, onde $A_{ij} = 1$ se existe uma aresta entre os
vértices $i$ e $j$, e $0$ caso contrário \parencite{newman2010}. Para
o exemplo acima:
\[
A = \begin{bmatrix}
0 & 1 & 1 & 0 \\
1 & 0 & 1 & 0 \\
1 & 1 & 0 & 1 \\
0 & 0 & 1 & 0
\end{bmatrix}
\]

\begin{figure}[ht]
    \centering
    \includegraphics[width=0.45\linewidth]{image2.png}
    \caption{Grafo exemplo (rede de colaboração entre investigadores).}
    \label{fig:grafo-exemplo}
\end{figure}

\subsection{Tipos de Grafos}
\label{sec:tipos-grafos}

Os grafos podem ser classificados de acordo com a \textbf{natureza
dos vértices e das arestas} \parencite{szwarcfiter2018}. Os principais
tipos relevantes para este trabalho são:

\begin{itemize}
    \item \textbf{Grafos não direcionados:} As arestas não têm
        orientação, ou seja, a relação é simétrica. No exemplo
        anterior, a colaboração é mútua e, portanto, o grafo é não
        direcionado.

    \item \textbf{Grafos direcionados (digrafos):} As arestas possuem
        uma direção, indicando uma relação assimétrica. Por exemplo,
        em uma rede de citações, se o artigo $A$ cita o artigo $B$,
        existe uma aresta de $A$ para $B$, mas não necessariamente o
        contrário. Formalmente, as arestas são pares ordenados
        $(u, v)$. A matriz de adjacência, nesse caso, não é simétrica.

    \item \textbf{Grafos ponderados:} As arestas possuem um peso
        associado, que pode representar intensidade, custo ou
        distância. Por exemplo, em uma rede de transportes, o peso
        pode ser a distância entre cidades. A matriz de adjacência
        armazena os pesos: $A_{ij} = w_{ij}$.

    \item \textbf{Grafos bipartidos:} Os vértices podem ser
        particionados em dois conjuntos disjuntos $U$ e $V$ de tal
        forma que todas as arestas conectam um vértice de $U$ a um
        vértice de $V$. Esse tipo de grafo é natural para modelar
        relações entre dois tipos distintos de entidades, como
        autores e artigos \parencite{newman2010}.

    \item \textbf{Grafos bipartidos ponderados:} Combinam as duas
        características anteriores: são grafos bipartidos cujas
        arestas possuem pesos. Um exemplo é uma rede de compras em
        que clientes (conjunto $U$) adquirem produtos (conjunto $V$)
        com determinada frequência (peso).
\end{itemize}

As classificações apresentadas seguem a sistematização clássica de
\citet{szwarcfiter2018} e \citet{newman2010}; a
Figura~\ref{fig:tipos-grafos} ilustra os cinco tipos discutidos.

\begin{figure}[ht]
    \centering
    \begin{subfigure}[t]{0.18\linewidth}
        \centering\includegraphics[width=\linewidth]{image3.png}
        \caption{Não direcionado}
    \end{subfigure}\hfill
    \begin{subfigure}[t]{0.18\linewidth}
        \centering\includegraphics[width=\linewidth]{image4.png}
        \caption{Direcionado}
    \end{subfigure}\hfill
    \begin{subfigure}[t]{0.18\linewidth}
        \centering\includegraphics[width=\linewidth]{image5.png}
        \caption{Ponderado}
    \end{subfigure}\hfill
    \begin{subfigure}[t]{0.18\linewidth}
        \centering\includegraphics[width=\linewidth]{image6.png}
        \caption{Bipartido}
    \end{subfigure}\hfill
    \begin{subfigure}[t]{0.18\linewidth}
        \centering\includegraphics[width=\linewidth]{image7.png}
        \caption{Bipartido ponderado}
    \end{subfigure}
    \caption{Tipos de grafos discutidos na
    Secção~\ref{sec:tipos-grafos}.}
    \label{fig:tipos-grafos}
\end{figure}

\subsection{Densidade: grafos esparsos e densos}
\label{sec:densidade}

Uma propriedade estrutural central de um grafo é a sua
\textbf{densidade}, que mede a fração de arestas efetivamente
presentes face ao número máximo possível \parencite{newman2010}. Para
um grafo \emph{direcionado} simples com $n = |V|$ vértices, o número
máximo de arestas é $n(n-1)$ --- cada par ordenado de vértices
distintos pode, quando muito, originar uma aresta ---, pelo que a
densidade se define como
\begin{equation}
    D = \frac{|E|}{n(n-1)} \in [0, 1].
    \label{eq:densidade}
\end{equation}
Para grafos não direcionados, o denominador é $n(n-1)/2$, dado que
cada aresta é não ordenada.

Diz-se que um grafo é \textbf{esparso} quando o número de arestas
cresce aproximadamente de forma linear com $n$ (densidade que tende
para zero à medida que $n$ aumenta) e \textbf{denso} quando $|E|$ se
aproxima do máximo $\Theta(n^2)$ (densidade próxima de $1$)
\parencite{newman2010, szwarcfiter2018}. A maioria dos grafos
oriundos de sistemas reais --- redes sociais, redes de citações ou
grafos de dependências de software --- é tipicamente \emph{esparsa}:
cada nó relaciona-se apenas com um pequeno subconjunto dos restantes.

Esta distinção é relevante para a presente dissertação por duas
razões. Em primeiro lugar, a densidade é usada como \emph{proxy} do
grau de \textbf{acoplamento arquitetural} de cada projeto analisado:
um Grafo de Requisitos mais denso indica responsabilidades mais
interligadas (ver Capítulo~\ref{ch:resultados},
Secção~\ref{sec:res-tec}, onde se reportam densidades de $0{,}125$,
$0{,}262$ e $0{,}304$ para os três casos de estudo, em ordem crescente
de acoplamento). Em segundo lugar, em grafos de pequena dimensão ---
como os aqui estudados, com poucas dezenas de vértices --- a densidade
é uma métrica \emph{volátil}: pequenas variações no número de arestas
deslocam-na de forma significativa, pelo que deve ser interpretada em
conjunto com outras métricas estruturais (presença de ciclos,
\emph{clustering}) e nunca isoladamente \parencite{newman2010}.

\subsection{Caminhos e passeios aleatórios}
\label{sec:caminhos}

Um caminho em um grafo é uma sequência de vértices
$v_1, v_2, \ldots, v_k$ tal que cada par consecutivo
$(v_i, v_{i+1})$ é uma aresta. O comprimento do caminho é o número de
arestas percorridas \parencite{newman2010}.

Um \textbf{passeio aleatório} (\emph{random walk}) é um processo
estocástico em que um ``caminhante'' se move ao longo do grafo de
forma probabilística \parencite{masuda2017, ross2014}. Em tempo
discreto, partindo de um vértice, o caminhante escolhe aleatoriamente
uma das arestas incidentes e move-se para o vértice vizinho. Esse
processo pode ser descrito por uma matriz de transição $P$, onde
$P_{ij}$ é a probabilidade de ir do vértice $j$ para o vértice $i$.
Para grafos não direcionados e não ponderados, uma escolha comum é
$P_{ij} = 1/\text{grau}(j)$ se $i$ e $j$ são vizinhos.

Passeios aleatórios são a base de diversos algoritmos de ranqueamento,
como o PageRank, que estima a importância de um vértice pela frequência
com que é visitado em um passeio aleatório de longa duração
\parencite{page1999, masuda2017}.

% ---------------------------------------------------------------------------
\section{Algoritmos de ranqueamento em grafos}
\label{sec:algoritmos}

O ranqueamento de vértices em grafos atribui um \emph{score} a cada
nó com base na estrutura de conexões. Diferentes algoritmos exploram
diferentes propriedades --- grau, centralidade de autovetor ou
passeios aleatórios --- e foram inicialmente propostos no contexto da
análise de redes de citações, do \emph{ranking} de páginas web e da
análise de redes sociais \parencite{page1999, kleinberg1999, lempel2001}.

\begin{notebox}[Importante]
    Todas as métricas descritas nesta secção são analisadas no
    contexto de \textbf{redes direcionadas}, dado que o grafo de
    requisitos é, por natureza, direcionado: a aresta
    $r_a \rightarrow r_b$ codifica a dependência de $r_a$ sobre $r_b$.
\end{notebox}

\paragraph{Escolha de PageRank, HITS e SALSA face a alternativas.}
A literatura de \emph{centrality measures} oferece um conjunto vasto
de alternativas: \emph{eigenvector centrality} \parencite{newman2010},
\emph{Katz centrality}, \emph{betweenness} e \emph{closeness}, ou
ainda decomposições $k$-\emph{core}. Esta dissertação restringe a
análise a PageRank, HITS e SALSA por três razões: (a) Silva et~al.\
\parencite{silva2018} estabelecem o PageRank como \emph{baseline} no
domínio --- comparar com algoritmos \emph{de mesma família} mantém o
estudo \emph{commensurate}; (b) HITS e SALSA introduzem a separação
\emph{hub}/\emph{authority}, essencial para distinguir requisitos
\emph{orquestradores} de requisitos \emph{críticos por dependência
direta}; (c) as três métricas são interpretáveis como passeios
aleatórios em cadeias de Markov, oferecendo uma fundamentação
estocástica comum \parencite{ding2002}. A extensão a outras métricas
(\eg\ \emph{Katz}, $k$-\emph{core}) é discutida como trabalho futuro
(Secção~\ref{sec:trabalho-futuro}).

\subsection{Centralidade de grau}
\label{sec:grau}

A centralidade de grau é a métrica mais simples de centralidade em
grafos \parencite{newman2010}. Para um vértice $v$, o
\textbf{grau de entrada} (em grafos direcionados) é o número de
arestas que chegam a ele, e o \textbf{grau de saída} é o número de
arestas que partem. Formalmente, para um grafo com matriz de
adjacência $A$, o grau de entrada de $v$ é $\sum_{u} A_{uv}$. Esta
métrica é a base do SALSA \parencite{lempel2001}, que na sua versão
básica equivale a uma ponderação dos graus.

\subsection{PageRank}
\label{sec:pagerank}

O \textbf{PageRank} (PR) foi proposto por \citet{page1999} e estima a
importância de um vértice como a \textbf{distribuição estacionária}
de um passeio aleatório em tempo discreto sobre o grafo
\parencite{brin1998, masuda2017}. A intuição é que um vértice é
importante se é apontado por outros vértices importantes ---
definição recursiva resolvida por iteração. O algoritmo incorpora um
fator de amortecimento $d$ (tipicamente $0{,}85$), que representa a
probabilidade de o caminhante continuar a seguir as arestas; com
probabilidade $1 - d$, o caminhante salta para um vértice aleatório
escolhido segundo uma distribuição de personalização (em geral,
\textbf{uniforme} --- todos os vértices têm idêntica probabilidade
$1/n$ de serem escolhidos). A equação fundamental é:
\begin{equation}
    \mathbf{r} = d\, M\, \mathbf{r} + (1 - d)\, \mathbf{e}
    \label{eq:pagerank}
\end{equation}
onde $\mathbf{r}$ é o vetor de \emph{ranks}, $M$ é a matriz de
transição estocástica (cada coluna soma 1) e $\mathbf{e}$ é o vetor
de personalização (uniforme, $e_i = 1/n$ para todo $i$). A matriz $M$
é construída a partir da matriz de adjacência: $M_{ij} = 1/\text{out}(j)$
se há uma aresta de $j$ para $i$, e $0$ caso contrário. Para vértices
sem arestas de saída (\emph{dead ends}), são feitos ajustes para
garantir a estocasticidade (Secção~\ref{sec:dangling}).

Trata-se de um \textbf{processo iterativo}: a equação
$\mathbf{r}^{(t+1)} = d M \mathbf{r}^{(t)} + (1-d)\mathbf{e}$ é
aplicada repetidamente a partir de uma estimativa inicial
$\mathbf{r}^{(0)}$. Sob as condições garantidas pelo amortecimento
(que torna a cadeia de Markov ergódica), a sequência
$\{\mathbf{r}^{(t)}\}_t$ converge para um único vetor estacionário
$\mathbf{r}^{*}$ \parencite{brin1998, masuda2017}. Na prática, a
iteração é interrompida quando
$\|\mathbf{r}^{(t+1)} - \mathbf{r}^{(t)}\|$ desce abaixo de uma
tolerância pré-definida.

\paragraph{Ilustração.} Considere o grafo direcionado simples com
vértices $A, B, C$ e arestas $A \to B$, $A \to C$, $B \to C$,
$C \to A$. Os graus de saída são $A : 2$, $B : 1$, $C : 1$, donde
\(
M =
\begin{psmallmatrix}
0 & 0 & 1 \\
0{,}5 & 0 & 0 \\
0{,}5 & 1 & 0
\end{psmallmatrix}
\)
(colunas somam $1$). Aplicando o PageRank com $d = 0{,}85$ e
$\mathbf{e}$ uniforme, obtém-se o vetor de \emph{ranks} após
convergência (Figura~\ref{fig:pagerank-exemplo}).

\begin{figure}[ht]
    \centering
    \includegraphics[width=0.4\linewidth]{image8.png}
    \caption{Grafo exemplo usado para ilustrar o PageRank.}
    \label{fig:pagerank-exemplo}
\end{figure}

\subsection{HITS (\emph{Hyperlink-Induced Topic Search})}
\label{sec:hits}

O \textbf{HITS}, proposto por \citet{kleinberg1999}, distingue dois
papéis para os vértices: \textbf{autoridades} (páginas com conteúdo
relevante) e \textbf{hubs} (páginas que apontam para boas
autoridades). A ideia é que uma boa autoridade é apontada por muitos
bons \emph{hubs}, e um bom \emph{hub} aponta para muitas boas
autoridades. Esta relação de reforço mútuo é expressa, de forma
iterativa, por:
\begin{equation}
    \mathbf{a} = A^{T}\mathbf{h}, \qquad \mathbf{h} = A\mathbf{a}
    \label{eq:hits}
\end{equation}
onde $A$ é a matriz de adjacência (para grafos direcionados). Após
normalização, os vetores de autoridade $\mathbf{a}$ e \emph{hub}
$\mathbf{h}$ convergem para os autovetores principais de $A^{T}A$ e
$AA^{T}$, respetivamente \parencite{kleinberg1999, ng2001}. A
Figura~\ref{fig:hits-exemplo} ilustra a aplicação ao mesmo grafo
exemplo.

\begin{figure}[ht]
    \centering
    \includegraphics[width=0.4\linewidth]{image9.png}
    \caption{Grafo exemplo usado para ilustrar o HITS.}
    \label{fig:hits-exemplo}
\end{figure}

\subsection{SALSA (\emph{Stochastic Approach for Link-Structure Analysis})}
\label{sec:salsa}

O \textbf{SALSA}, proposto por \citet{lempel2001}, combina ideias do
PageRank e do HITS num enquadramento estocástico. A sua intuição é a
seguinte: em vez de aplicar HITS diretamente sobre o grafo direcionado
original, o SALSA constrói \textbf{dois passeios aleatórios em um
grafo bipartido auxiliar} em que cada vértice do grafo original é
representado duas vezes --- uma como \emph{hub} e outra como
\emph{authority}. Os dois passeios alternam entre estes dois lados:

\begin{itemize}
    \item O passeio das \textbf{autoridades} parte de uma autoridade
        $v$, segue uma aresta no sentido inverso (até um \emph{hub}
        $u$ que a aponte) e depois segue uma aresta para a frente
        (até outra autoridade $w$ apontada por $u$). A distribuição
        estacionária deste passeio define o \emph{score} de
        autoridade de cada vértice.
    \item O passeio dos \textbf{hubs} é simétrico: parte de um
        \emph{hub}, segue uma aresta para a frente até uma autoridade
        e regressa, no sentido inverso, a outro \emph{hub}.
\end{itemize}

Sob normalização local (probabilidades $1/\text{out}(u)$ e
$1/\text{in}(v)$), em \textbf{grafos não ponderados e fortemente
conexos} a distribuição estacionária do passeio das autoridades
reduz-se a $\text{in}(v) / |E|$, ou seja, é diretamente proporcional
ao grau de entrada do vértice \parencite{lempel2001}. De forma
análoga, o \emph{score} de \emph{hub} é proporcional ao grau de
saída. Esta interpretação simples torna o SALSA computacionalmente
mais leve que o HITS, mas tem uma consequência empírica relevante
para esta dissertação: em grafos pequenos com graus pouco
diferenciados, o SALSA pode produzir \emph{distribuições praticamente
uniformes}, o que se traduz num \emph{ranking achatado}
(ver Capítulo~\ref{ch:resultados}, Secção~\ref{sec:discussao}).

\begin{figure}[ht]
    \centering
    \includegraphics[width=0.4\linewidth]{image10.png}
    \caption{Grafo exemplo usado para ilustrar o SALSA: os \emph{ranks}
    são proporcionais aos graus de entrada (\emph{authority}) e de
    saída (\emph{hub}).}
    \label{fig:salsa-exemplo}
\end{figure}

A vantagem teórica do SALSA face ao HITS é a sua resistência ao
\textbf{Efeito TKC} (\emph{Tightly Knit Community Effect}). Em grafos
onde existem \textbf{comunidades muito coesas} --- pequenos
subconjuntos de vértices densamente interligados --- o HITS tende a
concentrar quase toda a massa dos \emph{scores} nesses subgrupos,
ainda que sejam pouco representativos do grafo global. Isto acontece
porque o cálculo de autovetor reforça quem está fortemente
interligado, mesmo que isoladamente. O SALSA, ao normalizar
localmente por grau em cada passo do passeio, distribui a massa de
forma mais equilibrada e evita esta concentração espúria
\parencite{lempel2001}.

\begin{notebox}[Caveat metodológico]
    A vantagem do SALSA face ao Efeito TKC só se observa
    empiricamente quando o corpus contém grafos com sub-comunidades
    densas embutidas em estruturas maiores. Os três grafos analisados
    nesta dissertação são demasiado pequenos para apresentarem tal
    fenómeno; consequentemente, a propriedade só é confirmada no
    plano teórico e \textbf{não} é empiricamente observada (ver
    Capítulo~\ref{ch:resultados}, Secção~\ref{sec:discussao}).
\end{notebox}

\subsection{BiRank e normalizações simétricas}
\label{sec:birank}

O \textbf{BiRank}, proposto por \citet{he2014}, é um algoritmo
projetado para grafos bipartidos ponderados. A diferença em relação
ao PageRank reside na forma de \textbf{normalizar a matriz de pesos}
antes da iteração:

\begin{itemize}
    \item A \textbf{normalização estocástica} (usada pelo PageRank e
        por HITS) divide cada coluna $j$ da matriz $W$ pelo somatório
        dos pesos de saída do vértice $j$ --- ou seja, divide por
        $\sum_i W_{ij}$. Esta normalização torna as colunas
        estocásticas (somam $1$) e dá uma interpretação probabilística
        direta a cada transição.
    \item A \textbf{normalização simétrica} divide cada entrada
        $W_{ij}$ por $\sqrt{d_i \cdot d_j}$, onde $d_i$ e $d_j$ são
        os graus (ponderados) dos vértices envolvidos. O efeito desta
        segunda forma é \textbf{atenuar a influência desproporcional
        de vértices de grau muito elevado} (que dominariam, sob
        normalização estocástica), produzindo \emph{scores} mais
        estáveis em grafos com distribuições de grau muito
        assimétricas \parencite{he2014, zhou2004}.
\end{itemize}

A iteração do BiRank é:
\begin{equation}
    \mathbf{r}_{k+1} = \alpha S\mathbf{r}_{k} + (1 - \alpha)\mathbf{q}
    \label{eq:birank}
\end{equation}
onde $S = D_{u}^{-1/2} W D_{v}^{-1/2}$ é a matriz simetricamente
normalizada, $W$ é a matriz de pesos do grafo bipartido, $D_u$ e $D_v$
são as matrizes diagonais dos graus (em cada lado da bipartição) e
$\mathbf{q}$ é um vetor de consulta (\emph{prior information}). O
BiRank é flexível para extensões a grafos multipartidos
\parencite{he2015}.

\subsection{Comparação de rankings: correlação de Spearman e Kendall}
\label{sec:rank-correlation}

Os algoritmos descritos nas secções anteriores produzem, cada um, uma
\textbf{ordenação} dos vértices do grafo. Para avaliar em que medida
dois algoritmos concordam --- questão central desta dissertação
(Secção~\ref{sec:objetivos}) --- não basta comparar os \emph{scores}
numéricos brutos, que vivem em escalas diferentes; compara-se antes a
\textbf{ordem relativa} que cada algoritmo atribui aos vértices. Para
esse efeito recorre-se a dois coeficientes clássicos de
\textbf{correlação de ranks}, que constituem os \emph{critérios formais
de avaliação dos resultados} deste trabalho.

\paragraph{Coeficiente de Spearman ($\rho$).} O coeficiente de
Spearman \parencite{spearman1904} é a correlação de Pearson calculada
sobre os \emph{ranks} (posições ordinais) em vez dos valores
originais. Mede a força de uma relação \textbf{monotónica} entre dois
rankings: vale $+1$ quando as duas ordenações são idênticas, $-1$
quando são exatamente inversas e $0$ quando não há relação ordinal.
Na ausência de empates, admite a forma fechada
\[
    \rho = 1 - \frac{6 \sum_i d_i^2}{n(n^2 - 1)},
\]
onde $d_i$ é a diferença de posição do elemento $i$ entre os dois
rankings e $n$ o número de elementos.

\paragraph{Coeficiente de Kendall ($\tau_b$).} O coeficiente de
Kendall \parencite{kendall1938} baseia-se na contagem de
\textbf{pares concordantes e discordantes}: dois elementos formam um
par concordante se a sua ordem relativa é a mesma nos dois rankings, e
discordante caso contrário. A variante $\tau_b$
\parencite{goodman1954} introduz uma correção para \textbf{empates}
(\emph{ties}) --- frequentes em grafos pequenos, onde vários vértices
recebem \emph{scores} idênticos. Tende a ser mais conservador do que
$\rho$ e tem interpretação probabilística direta: aproxima a diferença
entre a probabilidade de concordância e a de discordância de um par
escolhido ao acaso.

Ambos os coeficientes tomam valores em $[-1, 1]$. A formalização
operacional usada na experimentação --- incluindo o tratamento de
\emph{ties} via \emph{mean ranks} e o critério \emph{a priori} de
``convergência forte'' ($\rho \geq 0{,}7$ \emph{e}
$\tau_b \geq 0{,}55$) --- é detalhada na Secção~\ref{sec:metricas}; os
valores obtidos para cada par de algoritmos em cada caso de estudo são
reportados e discutidos na Secção~\ref{sec:res-correlacoes}.

% ---------------------------------------------------------------------------
\section{Análise estática e Árvores de Sintaxe Abstrata (AST)}
\label{sec:ast}

Os algoritmos das secções anteriores operam sobre um grafo que, neste
trabalho, \textbf{não} é fornecido manualmente, mas \textbf{extraído
automaticamente do código-fonte}. O mecanismo que torna essa extração
possível é a \textbf{análise estática} via \emph{Abstract Syntax
Tree} (AST). Por fazer parte do fluxo principal do processo
(Capítulo~\ref{ch:metodologia}), o conceito é apresentado aqui.

\subsection{O que é uma AST}
\label{sec:ast-conceito}

A \textbf{análise estática} examina o código-fonte \emph{sem o
executar}, em contraste com a análise dinâmica, que observa o programa
em tempo de execução. O seu ponto de partida é a \textbf{Árvore de
Sintaxe Abstrata}: uma representação em árvore da estrutura sintática
do programa, em que cada nó corresponde a uma construção da linguagem
--- uma definição de função, uma chamada, uma atribuição, um laço.
Diz-se ``abstrata'' porque omite detalhes da escrita concreta
(parênteses, espaços, vírgulas) e retém apenas a estrutura
sintático-semântica relevante.

Em Python, a AST é produzida pelo módulo \code{ast} da biblioteca
padrão: a função \code{ast.parse(codigo)} devolve a árvore, que pode
ser percorrida para localizar nós de interesse. Para este trabalho, os
nós relevantes são as \textbf{definições de função}
(\code{FunctionDef}), os \textbf{métodos de classe} e, sobretudo, as
\textbf{invocações} (\code{Call}) --- é a partir destas que se infere
``que função chama que função''.

\paragraph{Exemplo.} Para o fragmento \code{def a(): return b()}, a
AST contém um nó \code{FunctionDef} (nome \code{a}) cujo corpo inclui
um nó \code{Call} ao identificador \code{b}. Percorrendo a árvore,
deteta-se a aresta de chamada $a \rightarrow b$ sem nunca executar o
código --- e, portanto, de forma \textbf{determinística}: a mesma
entrada produz sempre o mesmo grafo.

\subsection{Da AST ao Grafo de Requisitos: validação por fase}
\label{sec:ast-validacao}

A extração do Grafo de Requisitos a partir da AST não é um passo
único, mas um \emph{pipeline} de quatro fases ---
AST $\rightarrow$ \emph{Call Graph} $\rightarrow$ mapeamento
\code{func\_to\_req} $\rightarrow$ Grafo de Requisitos --- descrito em
detalhe na Secção~\ref{sec:reqgraph}. Coloca-se então uma questão
metodológica legítima: \emph{como se garante que o output de cada fase
está correto e que o grafo final representa fielmente as
dependências?} A resposta assenta em dois pilares.

\textbf{Determinismo e verificabilidade por fase.} Por se basear
exclusivamente em análise estática, cada fase produz um artefacto
\emph{inspecionável e reproduzível}, o que permite validar a correção
de forma \textbf{incremental} --- e não apenas no produto final:
\begin{itemize}
    \item A \textbf{AST} é gerada pelo \emph{parser} oficial da
        linguagem (módulo \code{ast}); a sua correção é garantida pela
        própria \emph{toolchain} do Python e não introduz
        variabilidade.
    \item O \textbf{Call Graph} é exportado como artefacto visual
        (\code{call\_graph.png}) e pode ser confrontado com a
        estrutura de chamadas conhecida do código --- arestas em falta
        ou espúrias tornam-se detetáveis por inspeção.
    \item O \textbf{Grafo de Requisitos} é exportado em formato
        legível (\code{req\_graph.json}, \code{req\_graph.dot}) e
        visual (\code{req\_graph.png}); a coerência das suas arestas
        com a arquitetura conhecida do projeto é verificada
        manualmente no passo de \emph{validação de coerência} do
        protocolo (Secção~\ref{sec:protocolo-func2req}).
\end{itemize}

\textbf{A noção de ``requisito'' adotada.} É importante sublinhar que
o conceito de requisito usado nesta dissertação \textbf{não} é o de
\citet{sommerville2007} --- um enunciado de necessidade do sistema ---,
mas sim uma \textbf{agregação de responsabilidades implementadas},
cujas interdependências são derivadas objetivamente da análise
estática do código (\cf\ nota ontológica na
Secção~\ref{sec:problema}). É precisamente esta opção que torna a
validação tratável: as dependências entre requisitos não são juízos
subjetivos a verificar contra a intenção dos \emph{stakeholders}, mas
\textbf{factos estruturais do código} --- se o requisito $r_a$ contém
uma função que invoca uma função do requisito $r_b$, a dependência
$r_a \rightarrow r_b$ existe \emph{por construção}. A correção do
grafo reduz-se, assim, à correção de duas coisas verificáveis: (i) a
deteção de chamadas pela AST (determinística) e (ii) o mapeamento
\code{func\_to\_req} --- o único passo com subjetividade residual,
cujo protocolo e limitações são tratados nas
Secções~\ref{sec:protocolo-func2req} e \ref{sec:ameacas}.

\begin{notebox}[Validação do grafo --- síntese]
    A validade do Grafo de Requisitos é \emph{construída fase a fase},
    não aferida apenas no fim: a AST é correta por garantia da
    \emph{toolchain}; o \emph{Call Graph} e o Grafo de Requisitos são
    artefactos inspecionáveis; e a semântica das arestas decorre
    objetivamente da estrutura do código. O único passo não
    verificável automaticamente --- o mapeamento \code{func\_to\_req}
    --- é isolado, com protocolo explícito
    (Secção~\ref{sec:protocolo-func2req}) e registado como ameaça à
    validade (Secção~\ref{sec:ameacas}).
\end{notebox}

% ---------------------------------------------------------------------------
\section{Aplicação das Métricas ao Ranking e Priorização de Requisitos}
\label{sec:re}

\subsection{Fundamentos da priorização de requisitos e desafios}
\label{sec:re-fundamentos}

A Engenharia de Requisitos (ER) é um domínio crucial da Engenharia de
Software, responsável pela definição e manutenção das necessidades do
sistema \parencite{sommerville2007}. Em projetos de software
complexos, a priorização de requisitos é uma atividade essencial para
o planeamento eficaz \parencite{karlsson2002, berander2005}.

A atividade de rastreabilidade de requisitos estabelece as relações
de dependência entre os requisitos, sendo essa a base estrutural para
a priorização \parencite{gotel1994, sommerville2007}. A literatura
identifica diversas técnicas tradicionais, como a \textbf{Comparação
em Pares} (\emph{Pair-wise Comparison}), que exige que os
\emph{stakeholders} comparem cada par de requisitos, sendo uma tarefa
demorada e de alto esforço \parencite{karlsson2007, karlsson1997}.
Outra técnica é a \textbf{Votação Acumulativa} (\emph{100-Points
Method}), que distribui pontos fixos entre \emph{stakeholders}, sendo
a prioridade final a soma dos pontos recebidos
\parencite{ahl2005, leffingwell2003}. Técnicas como o
\textbf{PERT/CPM} utilizam grafos para representar a sequência lógica
de planeamento, permitindo uma ordenação topológica
\parencite{kerzner2009}. O \textbf{AHP} (\emph{Analytic Hierarchy
Process}) \parencite{saaty1980} é, historicamente, o método mais
citado e usado em priorização de requisitos
\parencite{karlsson1997, berander2005}: estabelece uma hierarquia de
critérios e calcula pesos por comparação aos pares com verificação
de consistência, mas tem complexidade $O(n^2)$ que o torna inviável
para conjuntos com mais de algumas dezenas de requisitos.

Entretanto, métodos tradicionais de Priorização de Requisitos de
Software (PRS) sofrem com problemas intrínsecos de subjetividade e
esforço \parencite{firesmith2004, achimugu2014, bukhsh2020}. Fatores
como a inexperiência dos \emph{stakeholders}, a divergência na
interpretação das escalas de prioridade e o foco excessivo em apenas
um ponto de vista resultam em prioridades inconsistentes e ``posições
inválidas'' \parencite{silva2018}. Estudos práticos demonstram que a
Votação Acumulativa, por depender excessivamente da ação e da decisão
do \emph{stakeholder}, pode gerar um alto número de posições
inválidas e não resolve a duplicidade envolvida em requisitos
interdependentes \parencite{silva2018}. \emph{Surveys} recentes
\parencite{achimugu2014, bukhsh2020} confirmam estes desafios e
identificam o \emph{ranqueamento baseado em grafos} como uma das
direções mais promissoras para superar as limitações das técnicas
manuais.

\subsection{Construção do grafo de requisitos e posicionamento face
a Silva et~al.\ (2018)}
\label{sec:posicionamento}

A aplicação de algoritmos de ranqueamento estrutural na PRS baseia-se
na construção de um \textbf{grafo de requisitos}
\parencite{silva2018}. Neste modelo, cada requisito é um vértice e as
relações de dependência (obtidas pela rastreabilidade) são arestas
direcionadas. Especificamente, se o Requisito $A$ depende de $B$, um
\emph{link} de avanço origina-se em $A$ e aponta para $B$.

O presente trabalho posiciona-se explicitamente como uma
\textbf{evolução metodológica} do trabalho seminal de
\citet{silva2018}, que propõe a aplicação do PageRank à priorização
de requisitos. No estudo original, os autores: (i) restringem a
análise ao \textbf{PageRank} isoladamente, sem comparação sistemática
com outros algoritmos de ranqueamento; (ii) constroem o grafo de
requisitos a partir de \textbf{matrizes de rastreabilidade
preexistentes}, mantidas manualmente; e (iii) validam a abordagem com
uma avaliação qualitativa junto de profissionais. A presente
dissertação preserva o quadro conceptual de \citet{silva2018} ---
grafo de requisitos como suporte da priorização --- e estende-o em
três dimensões:

\begin{enumerate}
    \item \textbf{Alarga o leque de algoritmos}, comparando PageRank
        com HITS \parencite{kleinberg1999} e SALSA
        \parencite{lempel2001} sobre os mesmos grafos, recorrendo a
        métricas formais de \emph{rank correlation} (Spearman $\rho$,
        Kendall $\tau_b$) para quantificar a concordância.
    \item \textbf{Automatiza a construção do grafo de requisitos} a
        partir do código-fonte via análise estática (AST), eliminando
        a dependência de uma matriz de rastreabilidade mantida
        manualmente. Este \emph{trade-off} merece registo explícito:
        substitui-se uma fonte fiável mas dispendiosa (rastreabilidade
        validada por engenheiros) por outra mais escalável mas com
        validade de construto a discutir
        (Secção~\ref{sec:ameacas}, ameaças à validade).
    \item \textbf{Valida o pipeline sobre três projetos reais de
        código aberto} com perfis arquiteturais distintos, permitindo
        observar o comportamento dos algoritmos sob diferentes
        regimes de acoplamento.
\end{enumerate}

\subsection{Mapeamento dos algoritmos de ranqueamento para
priorização}
\label{sec:mapeamento}

Este referencial teórico fundamenta a aplicação de algoritmos de
ranqueamento de \emph{links} --- notavelmente PageRank, HITS e
SALSA --- para objetivar e reduzir o esforço na PRS
\parencite{silva2018, ding2002}. Esta abordagem transforma a rede de
dependências de requisitos num grafo estrutural, permitindo que a
priorização seja determinada pela topologia da rede, e não
predominantemente pela perceção subjetiva do \emph{stakeholder}.

O PageRank, em particular, tem-se mostrado sólido para a priorização
\parencite{silva2018}. Foi usado para medir a complexidade de
relacionamentos em sistemas \parencite{li2009} e para analisar o
impacto de preocupações em requisitos \parencite{jin2009}.

\subsection{Benefícios, limitações e considerações práticas}
\label{sec:beneficios}

A priorização via PageRank tem-se mostrado vantajosa em comparação
com métodos manuais como a Votação Acumulativa e o PERT/CPM,
oferecendo os seguintes benefícios \parencite{silva2018, he2014}:

\begin{enumerate}
    \item \textbf{Redução do esforço sem eliminação da subjetividade.}
        O PageRank diminui o envolvimento direto dos
        \emph{stakeholders}, focando-se na estrutura de dependência,
        e é escalável para lidar com um alto número de requisitos.
        Importa, contudo, sublinhar que a subjetividade é
        \emph{deslocada} (do julgamento direto para a construção do
        grafo), não eliminada --- ponto desenvolvido em
        Secção~\ref{sec:ameacas}.
    \item \textbf{Solução de Interdependência:} O PageRank, por meio
        das suas iterações, consegue solucionar a duplicidade
        envolvida em requisitos interdependentes, uma limitação não
        resolvida pela Votação Acumulativa ou PERT/CPM.
    \item \textbf{Ajuste Flexível de Prioridade:} É possível
        introduzir informações externas (prioridade técnica ou de
        negócio) no grafo por meio de \textbf{vértices artificiais}
        \parencite{silva2018}. Essa técnica opcional permite que o
        \emph{rank} de um requisito e o seu fecho transitivo de
        dependências sejam ajustados.
\end{enumerate}

Como consideração prática, embora o PageRank ainda possa resultar em
posições inválidas, esta abordagem estrutural revelou-se mais
favorável do que as técnicas manuais, apresentando menos
inconsistências e maior consistência com a lógica de dependência do
sistema \parencite{silva2018}. As limitações da abordagem incluem as
vulnerabilidades topológicas, como o \textbf{Efeito TKC}, que o SALSA
foi especificamente desenvolvido para mitigar \parencite{lempel2001}.
